\subsection{Independent survey: definitions and denominators}
We added a post hoc cross-sectional analysis of the 2022--2023 US National Survey of Family Growth (NSFG) public-use female respondent file \cite{nsfgdata}. This is a separate analysis, not a validation of Chinese GBD levels, China--G20 contrasts or GBD cause attribution. CDC's published estimates provided computational replication targets rather than new findings \cite{nugent2026}.

The original file contained 5,586 respondents. Our primary domain comprised 4,855 women aged 20--49 at interview. A sensitivity included one respondent aged 50 who had been 49 at screening according to the recode specification, yielding 4,856. Unknown diagnosis responses were excluded for each condition, not coded as no. Exposures were ever-reported professional diagnoses, not assigned causes of infertility.

Three estimands were distinguished: broad current fertility problems among all eligible women; that broad outcome among current opposite-sex married/cohabiting women; and 12-month infertility within the same union domain. Broad problems were FECUND categories 3--5 or INFERT=2. Surgical sterility (FECUND=1 or 2) was not classified as impaired fecundity. The union domain used RMARITAL categories 1--2, and the narrower outcome INFERT=2 with valid INFERT categories 1--3. The recodes can reflect either partner's fertility and do not isolate a disease-specific cause \cite{nugent2024,nsfgrecodes}.

Final weights WGT2022\_2023, strata VEST and clusters VECL were used. All 80 clusters in 20 strata were retained, including zero domain contributions, giving 60 design degrees of freedom. Taylor-linearized covariance across 36 proportions preserved dependence between overlapping definitions. Pointwise 95\% intervals used $t_{60}$, the log scale for ratios and the untransformed scale for differences. Paired changes in prevalence differences were evaluated in the same union domain. Analyses were exploratory, unadjusted and not multiple-comparison corrected. Narrow-outcome sparse-event counts were reported. R \texttt{survey} reconstructed proportions and the full covariance from original records using \texttt{svyratio} and a joint delta transformation of \texttt{svytotal}. The separate local Python calculation agreed to numerical precision; audit files record the comparison. This verifies computation under the design, not GBD validity.
